\documentclass[11pt]{article}
\usepackage[T1]{fontenc}
\usepackage[utf8]{inputenc}
\usepackage{lmodern,amsmath,amssymb,amsthm,booktabs,xurl}
\usepackage[margin=1in]{geometry}
\usepackage[colorlinks=true,urlcolor=blue,linkcolor=blue]{hyperref}
\newtheorem{proposition}{Proposition}
\newcommand{\lean}[1]{\nolinkurl{#1}}
\title{Local Bayesian Models, Global Coherence, and Revision}
\author{Mettapedia Working Note}
\date{November 2, 2026}
\begin{document}
\maketitle
\begin{abstract}
Bayesian revision requires a joint interpretation of the question and its
observations. This requirement does not entail a largest working context,
a materialized distribution over everything, or a uniquely determined coupling
of partially specified models. We distinguish these claims, identify the
extension assumptions under which local semantics admits an abstract global
measure, and connect the distinction to checked finite Bayesian conditioning.
Positive and negative examples delimit what local probability models license.
\end{abstract}

\section{Four different global claims}
A \emph{context} specifies quantities, their interpretations and a probability
model. Larger contexts may expose additional quantities. Four claims should
be kept separate:
\begin{enumerate}
\item A system stores or computes one universal joint distribution.
\item Its working contexts have a greatest member containing every quantity.
\item Its local laws admit at least one common joint interpretation.
\item Its retained information uniquely determines that interpretation.
\end{enumerate}
Claims 1 and 2 are unnecessary for local computation. Claim 3 is the substantive
coherence requirement when several contexts are asserted to describe the same
jointly meaningful classical quantities. Claim 4 need not hold: incomplete
knowledge can specify a family of admissible couplings.

For example, let contexts be finite subsets of a countably infinite variable
set. Two contexts have the common enlargement $C\cup D$. Every finite context
has a strict enlargement, so there is no greatest context. Nevertheless, an
infinite product of independent fair coins has exactly these coherent finite
marginals. Absence of a largest context and existence of a global measure are
therefore compatible. Conversely, two valid local models need not have a
compatible common enlargement when they impose additional overlapping
constraints.

\section{The precise revision law}
For a finite context with prior $p$, observation map $e$, query map $q$ and
nonnegative observation likelihood $\ell$, define
\[
 Z=\sum_x p(x)\ell(e(x)),\qquad
 p_E(x)=p(x)\ell(e(x))/Z \quad (Z>0).
\]
A zero normalizer does not define a posterior. Let $r(x)=(q(x),e(x))$.
\begin{proposition}[Context-independent evaluation on a fixed joint law]
The posterior query law depends only on $r_*p$, the joint law of the query
and observation. Any two finite contexts with the same $r_*p$ have the same
normalizer and, when it is positive, the same posterior query law.
\end{proposition}
\begin{proof}
The existing finite coarsening theorem gives
\[
 r_*\bigl(\operatorname{Bayes}(p,\ell\circ\pi_2\circ r)\bigr)
 =\operatorname{Bayes}(r_*p,\ell\circ\pi_2).
\]
Project the result along $\pi_1$. The same fibre-sum identity gives equality
of normalizers, so the zero-evidence boundary agrees as well.
\end{proof}
This is a theorem about fixed observations and a fixed interpreted likelihood.
Adding new evidence may change the answer. Retaining only a current posterior
also need not retain enough information for future evidence retraction.

A compatible directed family of finite contexts can use this proposition
without selecting a greatest context: each finite query and observation set
is evaluated in a suitable covering context, and compatibility makes the
answer independent of that choice. The existence of suitable covering
\emph{probability models} must be established, not inferred merely from the
existence of unions of variable names.

\section{Why separate marginals do not suffice}
Let $X$ and $Y$ both be fair bits. Three possible joint laws are:
\begin{center}
\begin{tabular}{lccc}
\toprule
Joint model & $P(X=1)$ & $P(Y=1)$ & $P(X=1\mid Y=1)$\\
\midrule
$X=Y$ & $1/2$ & $1/2$ & $1$\\
$X=1-Y$ & $1/2$ & $1/2$ & $0$\\
Independent & $1/2$ & $1/2$ & $1/2$\\
\bottomrule
\end{tabular}
\end{center}
Thus knowing the two marginals cannot determine how observing one quantity
revises the other. For disjoint variable sets, taking a product is a valid
coupling; choosing it asserts independence. A system may instead retain all
compatible couplings and return bounds. For two fair bits the conjunction
probability ranges from $0$ to $1/2$, and the displayed conditional ranges
from $0$ to $1$.

For two finite models $p(a,s)$ and $q(s,b)$ with the same overlap marginal
$r(s)$, a common enlargement can be constructed explicitly:
\[
 \mu(a,s,b)=\begin{cases}p(a,s)q(s,b)/r(s)&r(s)>0,\\0&r(s)=0.\end{cases}
\]
Summing over $b$ recovers $p$ and summing over $a$ recovers $q$; zero-marginal
rows vanish by nonnegativity. This construction selects conditional independence
of the private parts given the overlap. Compatibility permits this selection;
it does not establish that it matches the true dependence. More than two
contexts can impose incompatible cyclic constraints.

Even agreement on overlaps is insufficient. Consider three pair laws for
$(X,Y)$, $(Y,Z)$ and $(X,Z)$, each assigning $1/16$ to each equal outcome and
$7/16$ to each unequal outcome. Each pair is strictly positive and every
single-bit marginal is fair. But every Boolean triple has at least one equal
pair, hence any triple distribution must satisfy
\[
 P(X=Y)+P(Y=Z)+P(X=Z)\geq1.
\]
The specified pair laws give $3/8$, a contradiction. This is a classical
marginal-extension obstruction \cite{vorobev}; it is also the example in
\cite[Theorem 4.3]{omega}. It demonstrates that overlap consistency is weaker
than joint coherence. It does not establish that contradictory assessments
are simultaneously correct descriptions of one classical world. If quantities
are instead context-indexed, the cross-context identifications require their
own justification \cite{abramsky}.

\section{When an abstract global law follows}
Fix an index set $I$ and finite nonempty value sets $X_i$. Suppose a probability
law $\mu_C$ is supplied for \emph{every} finite $C\subseteq I$, with
\[
 (\operatorname{restrict}_{D,C})_*\mu_D=\mu_C\qquad(C\subseteq D).
\]
The Kolmogorov extension theorem yields a unique probability measure on
$\prod_{i\in I}X_i$ with its product sigma-algebra having these marginals.
One proof assigns measures to cylinder events using the finite laws;
compatibility makes this well-defined, and compactness of the product of finite
discrete spaces supplies the required extension property. Standard Borel
coordinate spaces admit the usual broader extension theorem. Arbitrary
measurable spaces require care.

Consequently, a system claiming mutually coherent exact classical probabilities
for all finite combinations over such a fixed vocabulary already commits to
an abstract global law in this sense. It can represent that law through local
models, kernels, programs or constraints. It need not construct a global table,
perform global inference on each update, or possess a greatest finite context.

This theorem does not apply to an arbitrary collection of pair marginals
without their higher-dimensional compatible laws. Nor does it supply a
canonical choice when only some marginals are known. Open-ended changes of
vocabulary or model class additionally require interpretation and compatibility
maps; finite working contexts alone do not settle those extension obligations.

\section{Probabilistic programs and truth-value interfaces}
A prior $\pi(m)$ over normalized probabilistic programs with common observable
semantics $p_m(x,e)$ gives a hierarchical model
\[
 p(m,x,e)=\pi(m)p_m(x,e),\qquad
 P(x\mid e)=\frac{\sum_m\pi(m)p_m(x,e)}{\sum_m\pi(m)p_m(e)}.
\]
Model weights revise by $\pi(m\mid e)\propto\pi(m)p_m(e)$. Averaging program
posteriors with unchanged prior weights generally gives a different answer.
Programs with different output meanings first need a common observation
interpretation. Weighted or nonterminating programs also require explicit
normalization or subprobability semantics.

This is a useful account of uncertainty over model structure. It does not
supply a unique joint distribution from separate truth-value summaries.
A Bernoulli parameter supplies a probability; a strength--confidence pair
also needs an evidence or parameter-uncertainty interpretation. A conditional
$P(B\mid A)$ differs from the material-implication probability
$P(\neg A\lor B)$, so a link's interpretation must be declared.

A transformer can supply a conditional distribution over token sequences.
Using it as an enriched truth-value object additionally requires an event
interpretation, model identity, conditioning information and an evidence
interface. Calibration, uncertainty over parameters and revision laws do not
follow from the presence of a tensor or a neural network.

\section{Specification assessment and implementation boundary}
MeTTaProN v0.3 already distinguishes proposal guidance from evidence admission,
requires observations to be scored at their sites, and says in Rule 10.3 that
lineage cannot supply statistical dependence. Section 17.4's packet admission
rule should explicitly require either a joint observation model or the
conditional-independence certificate that licenses multiplying packet
likelihoods. Disjoint provenance establishes non-reuse, not independence.
Successive \emph{conditional} likelihoods may also multiply by the chain rule;
they must then be interpreted relative to the recorded preceding observations.
A likelihood ratio alone does not determine an absolute evidence normalizer.

Section 17.5's compilation of conditional truth values is a model construction
requiring stated structural assumptions. Unrelated conditional summaries need
not specify a complete or consistent generative model. Merely multiplying
conditional factors on a cycle does not ensure that the normalized product
retains those factors as its conditionals.

The omega-PLN fidelity manuscript explicitly conditions its comparison on a
gluable covering context \cite[Section 2.2]{fidelity}. Its introductory phrase
``largest gluable context'' is stronger than required and needs an existence
argument. A specified compatible covering context suffices; an uncertain family
of completions can instead support lower and upper answers. If no completion
exists under the chosen identifications, the result is an incompatibility
report, not an ordinary global Bayesian answer.

\section{Checked status and remaining work}
The new Lean module
\lean{ProbabilityTheory/BayesianInference/LocalContexts.lean} reuses the existing
\lean{BayesianInference.coarsen_posterior} law from \lean{Coarsening.lean}.
It proves query--evidence reduction, agreement of normalizers and revised
answers between contexts, directedness without a greatest finite context,
the two-coupling revision control and the strictly positive three-pair
non-extension result. The theorem names include
\lean{query_from_joint}, \lean{covering_evidence_agree},
\lean{covering_contexts_agree}, \lean{finite_contexts_no_greatest},
\lean{Controls.coupling_revision} and \lean{Controls.no_joint_for_noisy_pairs}.

Existing indexed world-model and source-gluing modules supply context
infrastructure and further non-gluing controls. They do not themselves
establish probabilistic compatibility. The existing projective-credal layer
has a compact-carrier/finite-intersection theorem for compatible precise
previsions; its carrier, closedness and finite-feasibility assumptions remain
explicit. The sigma-additive extension statement above is classical background,
not a new theorem proved by LocalContexts. General infinite evidence and
continuous conditioning require their appropriate measure-theoretic hypotheses.

A concrete runtime integration should retain model and observation identities,
record the coupling or compatibility authority for each cross-context operation,
and distinguish a unique answer, bounded uncertainty, unavailable cross-context
semantics and impossible evidence. Those are interface obligations beyond the
finite mathematical result established here.

\begin{thebibliography}{9}
\bibitem{vorobev} N. N. Vorob'ev. Consistent Families of Measures and Their
Extensions. \emph{Theory of Probability and its Applications} 7(2), 147--163,
1962. \url{https://www.mathnet.ru/eng/tvp4710}.
\bibitem{abramsky} S. Abramsky and A. Brandenburger. The Sheaf-Theoretic
Structure of Non-Locality and Contextuality. \emph{New Journal of Physics} 13,
113036, 2011. \url{https://arxiv.org/abs/1102.0264}.
\bibitem{omega} B. Goertzel. Omega-PLN: Atlas-indexed reversible evidence-fibered
geodesic PLN, with local Bayesian charts, higher proof geometry, and forgetting
as holonomy. Manuscript v4, July 2026.
\bibitem{fidelity} B. Goertzel. Probably approximately Bayesian: Fidelity of
omega-PLN to global probabilistic reasoning, stability under approximate
implementation, bounded compounding of cached inference, and what a Hyperon
system must store, measure and learn. Manuscript v2, September 2026.
\bibitem{spec} B. Goertzel. MeTTaProN specification. Draft 0.3, 2026.
\end{thebibliography}
\end{document}
